% =====================================================================
% Evaluation: methodology + STRIDE security + operational cost + threats.
% Sprint 1 (consolidation): security reorganized with STRIDE (per [AB]); the
% four-architecture comparison (M0/A/B/SGX) is introduced and summarized in the
% comparative matrix (Table~\ref{tab:comparison}). Cells marked TBD there are
% measured/cited in the follow-up campaign (M0, boot breakdown, AWS cost, gas).
% Language: English. Tables: results_tables.tex + section_results_ab.tex.
% =====================================================================

The architectural argument of this paper is comparative: the compartmentalized
design should contain compromise better than less-separated alternatives while
paying an acceptable operational cost. We evaluate the implementation along two
axes. For \emph{security}, we organize the analysis with the STRIDE
model~\cite{shostack2014threat}, applied per interaction over the elements of
Figure~\ref{fig:architecture}, and show how the design neutralizes---or
explicitly fails to neutralize---each non-trivial threat, mapping every claim
to a concrete test. For \emph{cost}, we measure the latency, build
time, and resource price of the isolation strategy.

The comparison spans up to four points on an isolation spectrum, summarized in
Table~\ref{tab:comparison}: a no-isolation monolith (M0), a centralized
deployment that adds ephemeral microVMs but colocates all assets (Environment~A),
the compartmentalized design that also separates the planes (Environment~B), and,
as a cited reference, the SGX-based original DNAT~\cite{nascimento2020dnat}.
Because A and B share the same Firecracker microVMs, any cost or security
difference between them isolates the effect of \emph{compartmentalization}; M0
isolates the cost and the security delta of the \emph{microVM} primitive itself;
and the SGX reference frames the reuse and dependency limitations that motivate
our change of technology.

\subsection{Methodology}

We exercise the running prototype with a small suite of purpose-built
applications, each designed to attack one STRIDE category, over a synthetic,
PII-free dataset. Tests that target the execution plane submit an application
bundle directly to the executor, reproducing the same execution contract the
control plane generates, so that the variable under study---the execution
microVM---is isolated from unrelated services; tests that target lateral movement
inspect the executor environment, and the access-control test drives the
marketplace API. Each adverse application emits machine-readable evidence markers
that a harness collects, and each test is mapped to the interaction and STRIDE
category it supports (Table~\ref{tab:stride}), so that every observation
connects to a property of a specific boundary crossing rather than merely
listing successful checks.

Environments~A and~B were both measured live and end to end. The centralized
baseline (Environment~A) keeps execution under the same Firecracker isolation as
the compartmentalized design (Environment~B), so the only independent variable
between the two is the separation of planes; any measured difference therefore
reflects compartmentalization, not a weaker execution sandbox. The no-isolation
baseline (M0)---which runs the application directly in the control namespace,
without launching a microVM---was measured on the same host, and the on-chain gas
cost was measured on Hardhat; together these complete the comparative matrix
(Table~\ref{tab:comparison}), leaving only the SGX reference column, which is
cited from the original DNAT work.

\subsection{Security Evaluation with STRIDE}

STRIDE is an analysis of components and their interactions, not of a system
taken as a whole: a claim such as ``the design resists tampering'' is only
meaningful relative to a stated element. We therefore fix the level of
abstraction first. The analysis is performed over the elements of
Figure~\ref{fig:architecture}---the external parties, the three planes (each a
separate trust domain), the ledger, the store, and the package indexes---and
over the numbered interactions I1--I8 that cross their trust boundaries.
Following STRIDE-per-element practice~\cite{shostack2014threat}, each threat
category was examined against every element and boundary-crossing interaction;
flows internal to a single trust domain are not analyzed.
Table~\ref{tab:stride-matrix} records the full sweep---every category against
every interaction---classifying each cell as actively countered, trivially
handled, or an acknowledged residual, so that the coverage can be audited
rather than taken on trust. Table~\ref{tab:stride} then concentrates on the
non-trivial threats---those the architecture must actively counter---naming for
each the exact interaction where it arises, the countering mechanism, and the
test that provides evidence. The trivial and residual cells are discussed in
the text at the end of this subsection.

\begin{table}[htbp]
\centering
\caption{Completeness matrix: every STRIDE category against every interaction
of Figure~\ref{fig:architecture} (last two rows are element-level attack
surfaces, not legitimate flows). $\bullet$~countered by a design mechanism
(detailed in Table~\ref{tab:stride} or the text); $\circ$~trivial---standard
TLS/HTTPS or a trust assumption of Section~\ref{sec:proposal};
$\odot$~acknowledged residual or out of scope, discussed in the text and
Threats to Validity; ``--''~not applicable to that element or flow, or the
availability of external infrastructure (ledger, package indexes) assumed out
of scope. MicroVM/VMM escape at I7 is excluded by the trusted-VMM assumption of
Section~\ref{sec:proposal}; the reverse-pivot row captures what a plane
compromise reaches despite it.}
\label{tab:stride-matrix}
\footnotesize
\begin{tabular}{@{}lcccccc@{}}
\toprule
Interaction / element & S & T & R & I & D & E \\
\midrule
I1\quad external party $\leftrightarrow$ control & $\bullet$ & $\circ$ & $\bullet$ & $\circ$ & $\odot$ & $\circ$ \\
I2\quad control $\leftrightarrow$ ledger & $\circ$ & $\circ$ & $\bullet$ & $\circ$ & -- & -- \\
I3\quad control $\leftrightarrow$ store & $\circ$ & $\bullet$ & -- & $\bullet$ & -- & -- \\
I4\quad control $\rightarrow$ builder & $\odot$ & $\odot$ & -- & $\odot$ & -- & $\bullet$ \\
I5\quad builder $\rightarrow$ package index & $\circ$ & $\circ$ & -- & $\odot$ & -- & $\bullet$ \\
I6\quad builder $\rightarrow$ control & $\odot$ & $\bullet$ & -- & $\circ$ & -- & -- \\
I7\quad control $\rightarrow$ executor & $\odot$ & $\bullet$ & -- & $\bullet$ & -- & -- \\
I8\quad executor $\rightarrow$ control & $\odot$ & $\odot$ & -- & $\odot$ & -- & -- \\
\addlinespace
CVM3 $\rightarrow$ CVM1 (reverse pivot) & -- & -- & -- & -- & -- & $\bullet$ \\
CVM3 (per-run hygiene) & -- & -- & -- & -- & $\bullet$ & -- \\
\bottomrule
\end{tabular}
\end{table}

\begin{table}[htbp]
\centering
\caption{Non-trivial STRIDE threats, per interaction of
Figure~\ref{fig:architecture} (S/T/R/I/D/E = spoofing, tampering, repudiation,
information disclosure, denial of service, elevation of privilege). Each row
names the interaction where the threat arises, the countering mechanism, the
measured outcome, and the test providing it; trivial or assumption-covered
cells are discussed in the text.}
\label{tab:stride}
\scriptsize
\begin{tabularx}{\textwidth}{@{}p{1.15cm}cXXp{0.8cm}@{}}
\toprule
Interaction & Cat. & Threat \& countering mechanism & Outcome (measured) & Test \\
\midrule
I1 & S & Caller without access rights; on-chain gate (\texttt{hasAccessByIds}) before any bundle is prepared & Run rejected, \emph{access denied}, before code or data move & T5 \\
\addlinespace
I1, I2 & R & Denial of a registration or purchase; recorded as ledger transactions & Both actions recorded on chain, auditable & T5 \\
\addlinespace
I3 & T & Artifact modified between storage and use; SHA-256 content hash bound on chain & On-chain hash equals file SHA-256 & T5 \\
\addlinespace
I4, I5 & E & Malicious dependency runs at build time; ephemeral builder, filtered egress, no dataset & Payload runs in builder only---no secrets, no dataset; cache write denied & T2, T2b \\
\addlinespace
I6, I7 & T & Poisoned cache or modified artifact; mounted read-only at use & Build-time write denied (\texttt{Permission denied}) & T2b \\
\addlinespace
I7 & I & Dataset exfiltration over the network; execution microVM has no interface & 7/7 outbound channels fail, no \texttt{eth0} & T1, T2a \\
\addlinespace
I7 & I & Reads control-plane secrets or sibling data; no secrets, own disks only & Secrets absent from env; only own I/O disks visible & T1b, T4 \\
\addlinespace
CVM3\,$\to$\,CVM1 & E & Executor pivots into the control plane; secrets and stores absent & Keys, IPFS, cache, executions absent (\emph{residual:} shared net) & T4 \\
\addlinespace
CVM3 & D & Residual state starves later runs; per-run caps (1~vCPU/512~MiB) and teardown & No socket, disk, or mount after 9+ runs & T3 \\
\bottomrule
\end{tabularx}
\end{table}

\textbf{Reading the evidence.} Table~\ref{tab:stride} carries the
per-interaction outcomes; two points deserve emphasis beyond the data. First,
the containment of the execution plane is \emph{structural} rather than
policy-based: the plane that consumes the dataset has no network interface at
all, so the seven outbound probes fail with \texttt{Network is unreachable}
because there is no route to disable, not because a firewall rule denies
them---the strongest form the property can take. This is reinforced by code
size: the execution plane that runs untrusted code is also the smallest
first-party codebase (about 365 lines, against roughly 590 in the builder and
1{,}500 in the secret-holding control plane), so the boundary an attacker
reaches first is the least code to audit and holds none of the control-plane
logic. Second, the supply-chain test (T2) is the paper's thesis in miniature: a
package seeded into the builder's cache is resolved, bundled, and runs its
payload at import time inside the execution microVM---yet finds no secrets and a
blocked network, exactly mirroring T1 and T1b, while its build-time
\texttt{setup.py} variant executes in the isolated builder and fails against the
read-only cache. Malicious code, whether at build or execution time, inherits
the containment of the plane that runs it rather than reaching the control
plane.

\textbf{Trivial, assumption-covered, and residual cells.} The remaining cells
of the sweep in Table~\ref{tab:stride-matrix} fall into three groups.
\emph{Trivial.} On the external-facing interaction I1, tampering and
information disclosure in transit are handled by standard TLS. Interactions
with the ledger and the store (I2, I3) reduce to the trust assumptions of
Section~\ref{sec:proposal}: the ledger is assumed to execute correctly and the
store to be content-addressed, so spoofing and tampering there are outside the
design's control, and the access rights and content hashes they carry are
public by construction. One store cell is a mitigation we do claim rather than
assume: the artifacts are encrypted by the control plane before publication, so
information disclosure at the store (I3) is countered by that encryption.
Spoofing or tampering of a package index (I5) is the classic supply-chain
vector; the design does not authenticate the index beyond standard HTTPS and
instead contains what its packages can do once resolved (the I4/I5 row of
Table~\ref{tab:stride}). Two further trivial cells close the group: information
disclosure of the returned artifact (I6) is immaterial because the control plane
commissioned and already holds it, and elevation of privilege through the
marketplace API itself (I1) is a matter of contract and API implementation
correctness---an assumption outside the architectural analysis rather than a
property the planes enforce. \emph{Acknowledged residuals and out of scope.}
Three residuals are neither trivial nor fully closed in the prototype, and we
mark them as such rather than claim them. First, the internal service channels
(I4, I6--I8) are not mutually authenticated: server-side TLS does not by itself
stop a compromised plane from impersonating a legitimate internal caller, so
spoofing on these channels, in-transit tampering on the build bundle to the
builder (I4), and reach to services that should be unreachable are a
residual that internal mutual TLS and network segmentation would close, the same
network-reach residual already visible in Table~\ref{tab:blast-radius} and
revisited under Threats to Validity. Second, a
malicious dependency running in the builder could exfiltrate the buyer's own
application source over the filtered egress (I4/I5, information disclosure); the
damage is bounded because the provider's dataset never enters the build plane,
so only the buyer's own asset---not the sensitive data---is at risk, but we
record this as an accepted residual rather than a closed threat. Third, the
result returned to the buyer (I8) is out of scope in two senses: this work
bounds what a malicious application can \emph{reach}, not what an authorized
computation \emph{reveals} about its inputs (information disclosure), nor
whether a compromised executor returned a faithful result (tampering);
verifying the processing itself is a complementary line of work discussed in
Section~\ref{sec:related-work}. Denial of service beyond per-run state hygiene
(e.g., request flooding at I1) is likewise out of scope (Threats to Validity),
and repudiation is meaningful only for the ledger-recorded marketplace actions
already covered by the I1/I2 row.

Read against the requirements of Section~\ref{sec:problem-definition}, the
per-interaction analysis exercises R5 (rows I1 and I1/I2 and I3), R3 (rows
I4/I5 and I6), R1 and R4 (both I7 rows and the D row), and R2 (the
CVM3\,$\to$\,CVM1 row); R6 is not a threat-facing property and is exercised
instead by the cross-architecture comparison (Table~\ref{tab:comparison}).

\subsection{Operational Cost}

Table~\ref{tab:comparison} consolidates the cost of the isolation strategy
across architectures, and Table~\ref{tab:ab-performance} reports an earlier
live A/B campaign in detail. End-to-end execution latency is indistinguishable
between A and B; its absolute value tracks host warmth---26.8~s on the
2026-06-11 A/B campaign and 25.2~s on the warmer 2026-06-23 host used here for
the M0 comparison---but the A$\approx$B equality is invariant across both
campaigns. This time is dominated by
the microVM life cycle---kernel boot, disk discovery, output-disk synchronization,
and the guarded shutdown sequence---rather than by the application itself, whose
computation over sixty rows completes in milliseconds.

\begin{table}[htbp]
\centering
\caption{Comparative results across the isolation spectrum: a no-isolation
monolith (M0), a centralized deployment that adds ephemeral microVMs but
colocates all assets (A), the compartmentalized design that also separates the
planes (B), and the SGX-based original DNAT as a cited reference. Execution
latency and overhead are measured on the same host (2026-06-23); on-chain gas
is measured on Hardhat; build and artifact figures, identical across A and B,
are reported with the live A/B campaign (Table~\ref{tab:ab-performance}); the
SGX column is a qualitative reference (its enclave latencies are not
re-measured here).}
\label{tab:comparison}
\footnotesize
\begin{tabularx}{\textwidth}{@{}Xllll@{}}
\toprule
Metric & M0 & A & B & SGX (ref.) \\
\midrule
Execution latency, benign app (s)$^{\P}$ & 0.07 & 25.2 & 25.2 & n/a \\
microVM boot/life-cycle overhead (s) & $\approx$0 & $\approx$25.1 & $\approx$25.1 & n/a \\
On-chain gas, register / purchase (k)$^{\S}$ & 389 / 97 & 389 / 97 & 389 / 97 & --- \\
\addlinespace
App reuse with deps, no recompile & yes & yes & yes & \textbf{no} \\
\bottomrule
\end{tabularx}
\par\smallskip
{\footnotesize The microVM adds $\approx$25~s over M0's $\approx$0.07~s, since the
application's compute over 60 rows is sub-second; that fixed cost is the price of a
fresh, networkless microVM per run. A$\approx$B by design, as both share the same
Firecracker microVMs (build detail in Table~\ref{tab:ab-performance}, blast radius
in Table~\ref{tab:blast-radius}). $^{\P}$The execution figures are the clean-execution comparator (benign app)
from the 2026-06-23 same-host campaign; that campaign's three test apps spanned
25.2--25.8~s. Run-to-run dispersion is characterized in the earlier campaigns
(the $n{=}3$ A/B campaign at 26.8~s in Table~\ref{tab:ab-performance}, and a
colder $n{=}6$ campaign at 31.4~s, sd 1.4~s). On-chain gas is identical
across M0/A/B because the contract is the shared control plane; an asset's optional
bloom filter adds $\approx$185k gas per 256~B. $^{\S}$Gas is the EVM's unit of work
(a plain transfer is 21k); on a representative layer-2 (0.03~gwei, ETH at
US\$3,000) the cost is $\approx$US\$0.035 to register and $<$US\$0.01 to purchase.
The SGX column frames the reuse and dependency limitation of the original DNAT.}
\end{table}

The no-isolation baseline (M0) makes this concrete: running the same workload
directly in the container's namespace---no Firecracker, no disk isolation, no
network removal---the benign application completes in $\approx$0.07~s, so the
microVM adds roughly 25~s of fixed per-execution overhead. That overhead is the
price of launching a fresh, networkless microVM for every run, and it is exactly
what buys the containment: in the same measurement, M0 exposes a working
\texttt{eth0} and succeeds at network exfiltration, finds all six control-plane
secrets in its environment, and persists its writes, whereas A and B (with the
microVM) block exfiltration, expose no secret, and leave no residual state.
Separating the planes (B vs A) adds nothing on top of this microVM cost.

The on-chain cost is likewise architecture-independent. Registering an asset on
the marketplace contract costs $\approx$389k gas (a dataset with an empty bloom
filter; an application $\approx$327k), and purchasing access costs $\approx$97k
gas; a dataset's optional on-chain bloom filter adds $\approx$185k gas per 256
bytes. These figures are identical for M0, A, and B because the registry is the
shared control plane---isolation lives in the build and execution planes, not in
the contract---so the blockchain cost of the user journey is unchanged by the
isolation strategy. In concrete terms (a plain value transfer is 21k gas), on a
representative layer-2 deployment at 0.03~gwei with ETH at US\$3,000 these amount
to about US\$0.035 to register an asset and under US\$0.01 to purchase access, so
the marketplace is inexpensive to operate where such applications realistically
run.

The build plane is more expensive, as expected from its heavier configuration
(two vCPUs, network egress, and dependency resolution), but that cost is
identical across the two designs, which share the same builder microVM. The
per-path build latencies and artifact sizes are reported in detail with the
live A/B campaign in Section~\ref{sec:ab-results}
(Table~\ref{tab:ab-performance}).

\subsection{Threats to Validity}

The measurements were taken on a single host with a Linux-backed container
runtime; absolute latencies will vary with hardware (a colder first campaign saw
a 31.4~s execution mean and 103--131~s builds, with the same A$\approx$B
relationship), but the containment and exfiltration results are structural
(absence of a network interface, absence of secrets) and do not depend on timing.
The execution-plane tests bypass IPFS and the builder to isolate the microVM
under study; the end-to-end path through the control plane is exercised
separately by the access-control test. We report one residual weakness honestly:
in the single-host deployment the three planes share one Docker network, so a
compromised executor can still reach the control API and the blockchain RPC over
that network, even though the secrets and stores that define the blast radius
remain isolated; internal mutual TLS and network segmentation remain future work
(Table~\ref{tab:blast-radius}). Finally, the DoS analysis covers only per-run
state hygiene, not load-induced resource exhaustion, and the SGX column of
Table~\ref{tab:comparison} is qualitative---it restates the original DNAT's
reuse limitation rather than re-measuring enclave latencies here.

A more fundamental caveat concerns the execution substrate. The prototype runs
each plane as a container with Firecracker microVMs rather than inside a
hardware-backed confidential VM, so our measurements do not come from a genuinely
confidential environment and we make no claim about confidential-computing
overheads. This does not weaken the architectural argument: the properties we
evaluate---privilege separation across planes, a networkless execution microVM,
ephemeral teardown, and on-chain access control---are properties of the
architecture, not of the memory-encryption substrate. Running the same planes
inside confidential VMs (e.g., AMD SEV-SNP~\cite{amd2020sevsnp}) would change the
implementation and add a hardware-attestation and memory-encryption cost, but
would preserve the containment and blast-radius results demonstrated here. The
analysis therefore remains valid, with confidential execution as a substrate
substitution rather than an architectural change.
